% ===============================================================
\section{Что нужно помнить заранее}
% ===============================================================
Кратные интегралы опираются на несколько вещей из первых семестров. Перечислим их, чтобы дальше на них просто ссылаться.

\begin{enumerate}[label=\textbf{\arabic*.}]
\item \textbf{Определённый интеграл} $\int_a^b f(x)\,dx$ --- это площадь под графиком $y=f(x)$ на отрезке $[a,b]$ (со знаком: то, что ниже оси $Ox$, считается с минусом). Считается по формуле Ньютона--Лейбница
\[
\int_a^b f(x)\,dx=F(b)-F(a),\qquad F'(x)=f(x).
\]
\item \textbf{Замена переменной в определённом интеграле.} Если $x=g(t)$, то
\[
\int_{g(\alpha)}^{g(\beta)} f(x)\,dx=\int_\alpha^\beta f\bigl(g(t)\bigr)\,g'(t)\,dt .
\]
Обратите внимание: при замене (а) меняются пределы и (б) появляется множитель $g'(t)$ --- «во сколько раз растягивается кусочек $dt$». Двумерный и трёхмерный аналог этого множителя --- \emph{якобиан}; на нём держится вся тема.
\item \textbf{Интегрирование по частям:} $\displaystyle\int_a^b u\,dv=\bigl[uv\bigr]_a^b-\int_a^b v\,du$.
\item \textbf{Частная производная} $\dfrac{\partial f}{\partial x}$: дифференцируем по $x$, считая все остальные переменные константами. Например, для $f=x^2y^3$: $\ f'_x=2xy^3,\ f'_y=3x^2y^2$.
\item \textbf{Определители.}
$\begin{vmatrix} a&b\\ c&d\end{vmatrix}=ad-bc$; определитель $3\times 3$ раскладывается по строке (лучше --- по строке с нулями).
Определитель \emph{треугольной} матрицы (все элементы над или под диагональю --- нули) равен произведению диагональных элементов. И ещё: $\det(AB)=\det A\cdot\det B$.
\item \textbf{Тригонометрия:} $\sin^2t+\cos^2t=1$;\ \ $\cos^2t=\dfrac{1+\cos2t}{2}$;\ \ $\sin^2t=\dfrac{1-\cos2t}{2}$;\ \ $\sin t\cos t=\dfrac{\sin2t}{2}$;\ \ $1+\operatorname{tg}^2t=\dfrac1{\cos^2t}$.
\end{enumerate}

% ===============================================================
\section{Двойной интеграл}
% ===============================================================
\subsection{Откуда берётся двойной интеграл}

Пусть на плоскости $Oxy$ лежит область $\Omega$, а над ней --- «крыша» $z=f(x,y)\geqslant 0$. Тело, у которого дно --- $\Omega$, крыша --- график $f$, а боковые стенки вертикальны, называют \emph{цилиндрическим телом}. Как найти его объём?

\begin{center}
\begin{tikzpicture}[x={(-0.55cm,-0.4cm)},y={(1cm,0cm)},z={(0cm,1cm)},scale=1.1,>=Stealth]
  \draw[->] (0,0,0)--(3.6,0,0) node[below left]{$x$};
  \draw[->] (0,0,0)--(0,5,0) node[right]{$y$};
  \draw[->] (0,0,0)--(0,0,3.3) node[above]{$z$};
  % основание
  \fill[reg] plot[domain=0:360,samples=60,smooth] ({1.8+0.9*cos(\x)},{2.4+1.4*sin(\x)},0);
  \draw[regline] plot[domain=0:360,samples=60,smooth] ({1.8+0.9*cos(\x)},{2.4+1.4*sin(\x)},0);
  % крыша
  \draw[thick,regline] plot[domain=0:360,samples=60,smooth] ({1.8+0.9*cos(\x)},{2.4+1.4*sin(\x)},{2.1+0.3*sin(\x)-0.25*cos(\x)});
  % стенки
  \foreach \a in {60,180,250,320} \draw[gray] ({1.8+0.9*cos(\a)},{2.4+1.4*sin(\a)},0)--({1.8+0.9*cos(\a)},{2.4+1.4*sin(\a)},{2.1+0.3*sin(\a)-0.25*cos(\a)});
  % столбик
  \fill[orange!50] (1.9,2.3,0)--(1.9,2.6,0)--(1.9,2.6,2.15)--(1.9,2.3,2.15)--cycle;
  \draw[orange!70!black] (1.9,2.3,0)--(1.9,2.6,0)--(1.9,2.6,2.15)--(1.9,2.3,2.15)--cycle;
  \node at (2.5,2.0,0) {\small $\Omega$};
  \node[right] at (0,4.0,2.9) {\small $z=f(x,y)$};
  \draw[orange!60!black,->] (0,6.2,1.2) node[right]{\scriptsize объём столбика $\approx f(\xi_i,\eta_i)\,\Delta S_i$} -- (1.9,2.65,1.3);
\end{tikzpicture}
\end{center}

Разрежем $\Omega$ на маленькие кусочки $\Omega_1,\dots,\Omega_n$ с площадями $\Delta S_1,\dots,\Delta S_n$ и в каждом выберем точку $(\xi_i,\eta_i)$. Над маленьким кусочком тело почти не отличается от «столбика» высоты $f(\xi_i,\eta_i)$, поэтому
\[
V\approx\sum_{i=1}^n f(\xi_i,\eta_i)\,\Delta S_i .
\]
Чем мельче куски, тем точнее равенство. Это и есть идея двойного интеграла.

\begin{defbox}[\ (двойной интеграл)]
Сумма $\sum f(\xi_i,\eta_i)\,\Delta S_i$ называется \emph{интегральной суммой}. Пусть $\lambda$ --- наибольший из диаметров кусочков $\Omega_i$. Если при $\lambda\to0$ интегральные суммы стремятся к одному и тому же числу независимо от способа разбиения и выбора точек, это число называется \emph{двойным интегралом} функции $f$ по области $\Omega$:
\[
\iint\limits_{\Omega} f(x,y)\,dx\,dy=\lim_{\lambda\to0}\sum_{i=1}^n f(\xi_i,\eta_i)\,\Delta S_i .
\]
\end{defbox}

Запись $dx\,dy$ --- это «элемент площади» $dS$: если резать область прямыми, параллельными осям, то маленький кусочек --- прямоугольник со сторонами $dx$ и $dy$, его площадь $dx\,dy$.

\textbf{Смысл.} Если $f\geqslant0$ --- объём цилиндрического тела; если $f\equiv1$ --- площадь области: $\iint_\Omega dx\,dy=S(\Omega)$; если $f$ --- плотность пластинки --- её масса.

\begin{thmbox}[\ (существование)]
Если область $\Omega$ ограничена, замкнута, её граница состоит из конечного числа гладких кривых, а функция $f$ непрерывна на $\Omega$, то двойной интеграл существует. \emph{(Во всех наших задачах это выполнено.)}
\end{thmbox}

\textbf{Свойства} (те же, что у обычного интеграла):
\begin{enumerate}[label=\arabic*)]
\item линейность: $\iint(\alpha f+\beta g)=\alpha\iint f+\beta\iint g$;
\item аддитивность: если $\Omega$ разрезана на $\Omega_1$ и $\Omega_2$, то $\iint_\Omega=\iint_{\Omega_1}+\iint_{\Omega_2}$;
\item линии и точки имеют нулевую площадь, поэтому поведение функции (и замены переменных) на отдельных линиях на интеграл \emph{не влияет};
\item симметрия: если область симметрична относительно оси (плоскости) и функция при отражении не меняется, интеграл равен удвоенному интегралу по половине области.
\end{enumerate}

\subsection{Как считать: сведение к повторному интегралу}

Самое главное: двойной интеграл считают \emph{как два обычных}, один внутри другого.

Пусть область «правильная по $y$»: каждая вертикальная прямая $x=\mathrm{const}$ пересекает её по одному отрезку. Тогда область можно записать в виде
\[
\Omega=\{\,a\leqslant x\leqslant b,\quad y_1(x)\leqslant y\leqslant y_2(x)\,\}.
\]
Вертикаль «входит» в область через нижнюю кривую $y=y_1(x)$ и «выходит» через верхнюю $y=y_2(x)$.

\begin{center}
\begin{tikzpicture}[scale=1.0,>=Stealth]
  \draw[->] (-0.3,0)--(5.2,0) node[below]{$x$};
  \draw[->] (0,-0.3)--(0,3.6) node[left]{$y$};
  \fill[reg] plot[domain=1:4.2,smooth] (\x,{0.6+0.15*(\x-2.5)^2}) -- plot[domain=4.2:1,smooth] (\x,{3.0-0.12*(\x-2.2)^2}) -- cycle;
  \draw[thick,regline] plot[domain=1:4.2,smooth] (\x,{0.6+0.15*(\x-2.5)^2}) node[right]{\small $y=y_1(x)$};
  \draw[thick,regline] plot[domain=1:4.2,smooth] (\x,{3.0-0.12*(\x-2.2)^2}) node[right]{\small $y=y_2(x)$};
  \draw[regline] (1,{0.6+0.15*2.25})--(1,{3.0-0.12*1.44});
  \draw[regline] (4.2,{0.6+0.15*2.89})--(4.2,{3.0-0.12*4});
  \draw[dashed] (1,0) node[below]{$a$}--(1,0.94);
  \draw[dashed] (4.2,0) node[below]{$b$}--(4.2,1.03);
  \draw[->,very thick,hint] (2.8,0)--(2.8,{0.6+0.15*0.09}) node[pos=0.5,right]{\scriptsize вход};
  \draw[very thick,hint] (2.8,{0.6+0.15*0.09})--(2.8,{3.0-0.12*0.36});
  \draw[->,very thick,hint] (2.8,{3.0-0.12*0.36})--(2.8,3.5) node[right]{\scriptsize выход};
  \node[below] at (2.8,0) {$x$};
\end{tikzpicture}
\end{center}

\begin{thmbox}[\ (сведение двойного интеграла к повторному)]
Если $f$ непрерывна в области $\Omega=\{a\leqslant x\leqslant b,\ y_1(x)\leqslant y\leqslant y_2(x)\}$, то
\[
\iint\limits_\Omega f(x,y)\,dx\,dy=\int_a^b dx\int_{y_1(x)}^{y_2(x)} f(x,y)\,dy .
\]
Аналогично, если $\Omega=\{c\leqslant y\leqslant d,\ x_1(y)\leqslant x\leqslant x_2(y)\}$, то
\[
\iint\limits_\Omega f\,dx\,dy=\int_c^d dy\int_{x_1(y)}^{x_2(y)} f(x,y)\,dx .
\]
\end{thmbox}

\textbf{Как это читать.} Запись $\int_a^b dx\int_{y_1(x)}^{y_2(x)} f\,dy$ означает $\int_a^b\Bigl(\int_{y_1(x)}^{y_2(x)} f(x,y)\,dy\Bigr)dx$. Сначала считается \emph{внутренний} интеграл по $y$, причём $x$ в нём --- просто число. Получается функция от $x$, и её интегрируют по $x$.
Смысл: внутренний интеграл --- площадь сечения тела плоскостью $x=\mathrm{const}$; внешний «складывает» эти сечения.

\begin{notebox}[Три правила расстановки пределов]
\begin{enumerate}[label=\arabic*)]
\item Пределы \textbf{внешнего} интеграла --- всегда \textbf{числа}.
\item Пределы внутреннего интеграла могут зависеть только от внешних переменных.
\item Если нижняя (или верхняя) граница на разных участках задаётся разными формулами, область \textbf{режут на части} и пишут сумму интегралов (так будет в задаче 3).
\end{enumerate}
\end{notebox}

\textbf{Мини-пример.} $\Omega$ --- прямоугольник $0\leqslant x\leqslant1$, $0\leqslant y\leqslant2$. Тогда
$\iint_\Omega xy\,dx\,dy=\int_0^1 dx\int_0^2 xy\,dy=\int_0^1 x\cdot\frac{y^2}{2}\Big|_0^2dx=\int_0^1 2x\,dx=1.$
Для прямоугольника и функции-произведения $f=g(x)h(y)$ интеграл распадается в произведение: $\int_0^1x\,dx\cdot\int_0^2y\,dy=\frac12\cdot2=1$. Этим приёмом мы будем постоянно пользоваться: \emph{хорошая замена превращает область в прямоугольник, а подынтегральную функцию --- в произведение.}

\textbf{Смена порядка интегрирования} --- это та же область, но описанная «в другую сторону»: вместо вертикалей смотрим на горизонтали. Делается всегда по рисунку.

% ===============================================================
\section{Замена переменных в двойном интеграле}
% ===============================================================
\subsection{Зачем}
Если граница области --- окружности, гиперболы, лучи, параллелограмм, то в декартовых координатах пределы получаются неудобными (корни, разбиение на куски). Идея: ввести новые переменные $(u,v)$,
\[
x=x(u,v),\qquad y=y(u,v),
\]
так, чтобы граница превратилась в линии $u=\mathrm{const}$, $v=\mathrm{const}$. Тогда новая область --- прямоугольник, и все пределы --- числа.

Линии $u=\mathrm{const}$ и $v=\mathrm{const}$ рисуют на плоскости $Oxy$ «криволинейную сетку». Например, для полярных координат это окружности $r=\mathrm{const}$ и лучи $\varphi=\mathrm{const}$.

\subsection{Якобиан --- во сколько раз растягивается площадь}

Возьмём в плоскости $(u,v)$ маленький прямоугольник со сторонами $du$, $dv$. Куда он перейдёт?

\begin{center}
\begin{tikzpicture}[>=Stealth]
  \begin{scope}
  \draw[->] (-0.2,0)--(3,0) node[below]{$u$};
  \draw[->] (0,-0.2)--(0,2.6) node[left]{$v$};
  \fill[newreg] (1,0.8) rectangle (2,1.6);
  \draw[thick,newline] (1,0.8) rectangle (2,1.6);
  \node[below] at (1.5,0.8) {\scriptsize $du$};
  \node[right] at (2,1.2) {\scriptsize $dv$};
  \node[below left] at (1,0.8) {\scriptsize $(u,v)$};
  \end{scope}
  \draw[->,thick] (3.4,1.2) to[bend left=20] node[above]{\small $x=x(u,v)$} node[below]{\small $y=y(u,v)$} (5.4,1.2);
  \begin{scope}[xshift=6cm]
  \draw[->] (-0.2,0)--(3.6,0) node[below]{$x$};
  \draw[->] (0,-0.2)--(0,2.6) node[left]{$y$};
  \fill[reg] (1,0.6) to[bend right=8] (2.4,1.0) to[bend left=5] (2.7,2.0) to[bend left=8] (1.3,1.6) to[bend right=5] cycle;
  \draw[thick,regline] (1,0.6) to[bend right=8] (2.4,1.0) to[bend left=5] (2.7,2.0) to[bend left=8] (1.3,1.6) to[bend right=5] cycle;
  \draw[->,very thick,hint] (1,0.6)--(2.4,1.0) node[pos=0.75,below=3pt]{\scriptsize $\vec a$};
  \draw[->,very thick,hint] (1,0.6)--(1.3,1.6) node[midway,left]{\scriptsize $\vec b$};
  \node[left] at (1,0.6) {\scriptsize $M$};
  \end{scope}
\end{tikzpicture}
\end{center}

Вершина $(u,v)$ переходит в точку $M=(x(u,v),y(u,v))$. Соседняя вершина $(u+du,v)$ --- в точку, сдвинутую на вектор
\[
\vec a=\bigl(x(u+du,v)-x(u,v),\ y(u+du,v)-y(u,v)\bigr)\approx\Bigl(\frac{\partial x}{\partial u}du,\ \frac{\partial y}{\partial u}du\Bigr)
\]
(по определению производной). Аналогично вершина $(u,v+dv)$ сдвигается на $\vec b\approx\bigl(\frac{\partial x}{\partial v}dv,\ \frac{\partial y}{\partial v}dv\bigr)$.
Значит, образ маленького прямоугольника --- почти параллелограмм, натянутый на $\vec a,\vec b$. Площадь параллелограмма, натянутого на векторы $(a_1,a_2)$ и $(b_1,b_2)$, равна $|a_1b_2-a_2b_1|$, то есть модулю определителя. Получаем
\[
dS_{xy}\approx\left|\det\begin{pmatrix}\dfrac{\partial x}{\partial u}&\dfrac{\partial x}{\partial v}\\[2mm]\dfrac{\partial y}{\partial u}&\dfrac{\partial y}{\partial v}\end{pmatrix}\right|du\,dv .
\]

\begin{defbox}[\ (якобиан)]
Определитель
$\displaystyle J=\frac{\partial(x,y)}{\partial(u,v)}=\begin{vmatrix}x'_u&x'_v\\ y'_u&y'_v\end{vmatrix}$
называется \emph{якобианом} замены. Его модуль $|J|$ показывает, во сколько раз замена растягивает площади вблизи данной точки: $dx\,dy=|J|\,du\,dv$.
\end{defbox}

\begin{thmbox}[\ (замена переменных в двойном интеграле)]
Пусть
\begin{enumerate}[label=\arabic*)]
\item формулы $x=x(u,v)$, $y=y(u,v)$ взаимно однозначно переводят область $\Omega_{uv}$ в область $\Omega_{xy}$;
\item функции $x(u,v),y(u,v)$ имеют непрерывные частные производные;
\item якобиан $J\neq0$ внутри $\Omega_{uv}$;
\item $f$ непрерывна на $\Omega_{xy}$.
\end{enumerate}
Тогда
\[
\boxed{\ \iint\limits_{\Omega_{xy}} f(x,y)\,dx\,dy=\iint\limits_{\Omega_{uv}} f\bigl(x(u,v),y(u,v)\bigr)\,|J(u,v)|\,du\,dv.\ }
\]
Условия 1) и 3) разрешается нарушать на множестве нулевой площади (на отдельных точках и линиях, например на границе).
\end{thmbox}

Пример такого нарушения: в полярных координатах вся линия $r=0$ переходит в одну точку $O$ и там $J=0$ --- формула всё равно верна, потому что линия имеет нулевую площадь.

\textbf{Почему модуль?} В двойном интеграле пределы всегда ставятся от меньшего к большему, а площадь положительна. Знак якобиана говорит лишь о том, «перевернулась» ли ориентация --- для площади это неважно.

\subsection{Как найти новую область (график 2)}
\begin{enumerate}[label=\arabic*)]
\item Каждую линию границы (1), (2), \dots\ переписываем в новых переменных.
\item Определяем, по какую сторону от каждой линии лежит область (переписываем неравенства, задающие область, или подставляем пробную точку).
\item Рисуем новую область. Хорошая замена делает её прямоугольником (или хотя бы областью, где пределы простые).
\end{enumerate}

\subsection{Два полезных свойства якобиана}
\textbf{(а) Обратная замена.} Часто новые переменные удобно задать формулами $u=u(x,y)$, $v=v(x,y)$ (как в задачах 1, 5, 8). Тогда не обязательно выражать $x,y$ через $u,v$:
\[
\frac{\partial(x,y)}{\partial(u,v)}=\frac{1}{\ \dfrac{\partial(u,v)}{\partial(x,y)}\ } .
\]
Почему: если сделать замену и сразу обратную, получится тождественное отображение, его матрица производных --- единичная. По правилу дифференцирования сложной функции матрицы производных перемножаются, а $\det(AB)=\det A\det B$, значит, произведение двух якобианов равно $1$.

\textbf{(б) Замена в два этапа.} Если $(u,v)\to(p,q)\to(x,y)$, то по той же причине якобиан всей замены равен \emph{произведению} якобианов этапов. Так удобно считать якобиан сферических координат (см.\ ниже).

\subsection{Полярные координаты}
\begin{minipage}{0.55\linewidth}
\[
\begin{cases}x=r\cos\varphi,\\ y=r\sin\varphi,\end{cases}\qquad r\geqslant0,\ \ 0\leqslant\varphi<2\pi .
\]
$r=\sqrt{x^2+y^2}$ --- расстояние от точки до начала координат, $\varphi$ --- угол от положительной полуоси $Ox$ (против часовой стрелки).
\[
J=\begin{vmatrix}\cos\varphi&-r\sin\varphi\\ \sin\varphi&r\cos\varphi\end{vmatrix}=r\cos^2\varphi+r\sin^2\varphi=r,
\]
\[|J|=r .
\]
Геометрически: «прямоугольник» $dr\times d\varphi$ --- это кусочек кольца со сторонами $dr$ и $r\,d\varphi$ (длина дуги), его площадь $r\,dr\,d\varphi$.
\end{minipage}\hfill
\begin{minipage}{0.4\linewidth}\centering
\begin{tikzpicture}[scale=1.1,>=Stealth]
  \draw[->] (-0.3,0)--(3.2,0) node[below]{$x$};
  \draw[->] (0,-0.3)--(0,3.0) node[left]{$y$};
  \fill[reg] (35:1.8) arc(35:50:1.8) -- (50:2.4) arc(50:35:2.4) -- cycle;
  \draw[regline,thick] (35:1.8) arc(35:50:1.8) -- (50:2.4) arc(50:35:2.4) -- cycle;
  \draw[dashed] (0,0)--(35:2.9);
  \draw[dashed] (0,0)--(50:2.9);
  \draw (0.6,0) arc(0:35:0.6); \node at (17:0.85) {\scriptsize $\varphi$};
  \node[right] at (30:2.15) {\scriptsize $dr$};
  \node[above] at (45:2.5) {\scriptsize $r\,d\varphi$};
  \draw[gray] (0,0) -- (42:1.8) node[pos=0.5,above left]{\scriptsize $r$};
\end{tikzpicture}
\end{minipage}

\textbf{Словарик полярных координат:}
\begin{center}
\begin{tabular}{ll}
\toprule
в декартовых & в полярных\\
\midrule
окружность $x^2+y^2=R^2$ & $r=R$\\
луч $y=kx$ ($x>0$) & $\varphi=\operatorname{arctg}k$\\
прямая $x=c$ & $r=c/\cos\varphi$\\
прямая $y=c$ & $r=c/\sin\varphi$\\
функция от $\sqrt{x^2+y^2}$ & функция от $r$\\
\bottomrule
\end{tabular}
\end{center}
\textbf{Как ставить пределы:} обычно внешний --- угол $\varphi$ (между какими лучами зажата область), внутренний --- $r$: идём по лучу из начала координат; где вошли в область --- нижний предел, где вышли --- верхний.

\textbf{Обобщённые полярные координаты.} $x=a\,r\cos^n\varphi$, $y=b\,r\sin^n\varphi$ ($0\leqslant\varphi\leqslant\frac\pi2$). Якобиан:
\[
\begin{aligned}J&=\begin{vmatrix}a\cos^n\varphi&-na\,r\cos^{n-1}\varphi\sin\varphi\\ b\sin^n\varphi&nb\,r\sin^{n-1}\varphi\cos\varphi\end{vmatrix}
\\&=nab\,r\cos^{n-1}\varphi\sin^{n-1}\varphi\,(\cos^2\varphi+\sin^2\varphi)=nab\,r\cos^{n-1}\varphi\sin^{n-1}\varphi .\end{aligned}
\]
При $n=1$, $a=b=1$ получаются обычные полярные координаты.

\subsection{Линейная замена}
$u=a_1x+b_1y$, $v=a_2x+b_2y$. Тогда $\dfrac{\partial(u,v)}{\partial(x,y)}=a_1b_2-a_2b_1$ --- число, и $|J|=\dfrac{1}{|a_1b_2-a_2b_1|}$. Такая замена переводит параллелограмм со сторонами $u=\mathrm{const}$, $v=\mathrm{const}$ в прямоугольник (задача 1).

% ===============================================================
\section{Тройной интеграл}
% ===============================================================
Определение дословно повторяет двумерное: тело $V$ режется на маленькие кусочки объёма $\Delta V_i$,
\[
\iiint\limits_V f(x,y,z)\,dx\,dy\,dz=\lim_{\lambda\to0}\sum_i f(\xi_i,\eta_i,\zeta_i)\,\Delta V_i .
\]
Смысл: при $f\equiv1$ --- \textbf{объём тела} $V=\iiint_V dx\,dy\,dz$; если $f$ --- плотность, то масса.

\subsection{Сведение к повторному: «проекция + отрезок»}
Пусть каждая вертикальная прямая пересекает тело по одному отрезку: снизу тело ограничено поверхностью $z=z_1(x,y)$ («пол»), сверху --- $z=z_2(x,y)$ («крыша»), а $D$ --- проекция тела на плоскость $Oxy$. Тогда
\[
\boxed{\ \iiint\limits_V f\,dx\,dy\,dz=\iint\limits_D dx\,dy\int_{z_1(x,y)}^{z_2(x,y)} f(x,y,z)\,dz,\qquad
\text{объём:}\ \ V=\iint\limits_D\bigl(z_2(x,y)-z_1(x,y)\bigr)\,dx\,dy.\ }
\]

\begin{notebox}[Как найти проекцию $D$]
\begin{itemize}
\item Поверхности, в уравнения которых \textbf{не входит $z$} (например $y=x$, $y=x^2$, $x+y=1$), --- это вертикальные «стенки». На плоскость $Oxy$ они проектируются в линии с тем же уравнением.
\item Линия пересечения пола и крыши: приравниваем $z_1(x,y)=z_2(x,y)$ --- получаем уравнение без $z$, это тоже кусок границы $D$.
\item Чтобы понять, какая поверхность --- пол, а какая --- крыша, подставляем пробную точку из $D$.
\end{itemize}
\end{notebox}

Иногда удобнее \textbf{метод сечений}: если сечение тела плоскостью $z=\mathrm{const}$ --- простая фигура $D(z)$, то
$\iiint_V f=\int_{z_{\min}}^{z_{\max}}dz\iint_{D(z)}f\,dx\,dy$ (используем в задаче 8 для проверки).

% ===============================================================
\section{Замена переменных в тройном интеграле}
% ===============================================================
Всё как в плоскости, только матрица $3\times3$:
\[
J=\frac{\partial(x,y,z)}{\partial(u,v,w)}=\begin{vmatrix}x'_u&x'_v&x'_w\\ y'_u&y'_v&y'_w\\ z'_u&z'_v&z'_w\end{vmatrix},\qquad
\iiint\limits_{V} f\,dx\,dy\,dz=\iiint\limits_{V_{uvw}} f\bigl(x(u,v,w),\dots\bigr)\,|J|\,du\,dv\,dw .
\]
$|J|$ --- коэффициент растяжения объёма: маленький параллелепипед $du\,dv\,dw$ переходит в почти параллелепипед, натянутый на векторы $(x'_u,y'_u,z'_u)du$, $(x'_v,\dots)dv$, $(x'_w,\dots)dw$, а объём параллелепипеда равен модулю определителя из координат векторов. Условия теоремы --- те же, что в двумерном случае (нарушения допускаются на множествах нулевого объёма: точках, линиях, поверхностях).

\subsection{Цилиндрические координаты}
\[
\begin{cases}x=r\cos\varphi,\\ y=r\sin\varphi,\\ z=z,\end{cases}\qquad |J|=r .
\]
(Определитель $3\times3$ раскладывается по последней строке $(0,0,1)$ и сводится к полярному.) Удобны, когда в уравнениях встречается $x^2+y^2$:
параболоид $z=x^2+y^2\ \to\ z=r^2$; конус $z=\sqrt{x^2+y^2}\ \to\ z=r$; цилиндр $x^2+y^2=R^2\ \to\ r=R$.

\subsection{Сферические координаты}
\begin{minipage}{0.56\linewidth}
\[
\begin{cases}x=r\cos\psi\cos\varphi,\\ y=r\cos\psi\sin\varphi,\\ z=r\sin\psi,\end{cases}\qquad
\begin{array}{l} r\geqslant0,\\ 0\leqslant\varphi<2\pi,\\ -\frac\pi2\leqslant\psi\leqslant\frac\pi2.\end{array}
\]
$r=\sqrt{x^2+y^2+z^2}$ --- расстояние до начала координат, $\varphi$ --- «долгота» (угол в плоскости $Oxy$ от оси $Ox$), $\psi$ --- «широта» (угол между радиус-вектором и плоскостью $Oxy$; $\psi>0$ --- верхнее полупространство).
\end{minipage}\hfill
\begin{minipage}{0.4\linewidth}\centering
\begin{tikzpicture}[x={(-0.55cm,-0.45cm)},y={(1cm,0cm)},z={(0cm,1cm)},scale=1.2,>=Stealth]
  \draw[->] (0,0,0)--(2.4,0,0) node[below left]{$x$};
  \draw[->] (0,0,0)--(0,2.6,0) node[right]{$y$};
  \draw[->] (0,0,0)--(0,0,2.4) node[above]{$z$};
  \coordinate (P) at ({1.8*cos(35)*cos(50)},{1.8*cos(35)*sin(50)},{1.8*sin(35)});
  \coordinate (Q) at ({1.8*cos(35)*cos(50)},{1.8*cos(35)*sin(50)},0);
  \draw[thick,hint] (0,0,0)--(P) node[above right]{\small $M$};
  \draw[dashed] (0,0,0)--(Q)--(P);
  \draw (0.7,0,0) arc[start angle=0,end angle=50,radius=0.7] ;
  \node at ({0.95*cos(25)},{0.95*sin(25)},0) {\scriptsize $\varphi$};
  \node at ({1.1*cos(18)*cos(50)},{1.1*cos(18)*sin(50)},{1.1*sin(18)}) {\scriptsize $\psi$};
  \node at ({0.9*cos(35)*cos(50)},{0.9*cos(35)*sin(50)},{0.9*sin(35)+0.25}) {\scriptsize $r$};
\end{tikzpicture}
\end{minipage}

\textbf{Якобиан --- в два этапа} (свойство (б)). Этап 1: цилиндрические координаты $x=\rho\cos\varphi$, $y=\rho\sin\varphi$, $z=z$, якобиан $\rho$. Этап 2: в вертикальной полуплоскости $(\rho,z)$ вводим «полярные» координаты $\rho=r\cos\psi$, $z=r\sin\psi$, якобиан $r$. Итого
\[
\boxed{\ |J|=\rho\cdot r=r^2\cos\psi .\ }
\]
\emph{Замечание.} В части учебников вместо широты берут угол $\theta$ от оси $Oz$ ($z=r\cos\theta$), тогда $|J|=r^2\sin\theta$ --- это то же самое, так как $\theta=\frac\pi2-\psi$.

\textbf{Словарик:} сфера $x^2+y^2+z^2=R^2\to r=R$; сфера $x^2+y^2+z^2=2az\to r=2a\sin\psi$; конус $x^2+y^2=z^2$ ($z\geqslant0$) $\to\psi=\frac\pi4$; плоскость $z=0\to\psi=0$; плоскости $y=0$, $x=0$ (при $x,y\geqslant0$) $\to\varphi=0$, $\varphi=\frac\pi2$.

\subsection{Обобщённые сферические координаты}
Для тел вида $\left(\frac xa\right)^{2/n}+\left(\frac yb\right)^{2/n}+\left(\frac zc\right)^{2/n}\leqslant1$ (в первом октанте) берут
\[
\begin{cases}x=a\,r\cos^n\varphi\cos^n\psi,\\ y=b\,r\sin^n\varphi\cos^n\psi,\\ z=c\,r\sin^n\psi,\end{cases}\qquad 0\leqslant\varphi,\psi\leqslant\tfrac\pi2 .
\]
Тогда $\left(\frac xa\right)^{2/n}+\left(\frac yb\right)^{2/n}+\left(\frac zc\right)^{2/n}=r^{2/n}\bigl(\cos^2\varphi\cos^2\psi+\sin^2\varphi\cos^2\psi+\sin^2\psi\bigr)=r^{2/n}$, и поверхность превращается в $r=1$.

Якобиан снова в два этапа. Этап 1: $x=a\rho\cos^n\varphi$, $y=b\rho\sin^n\varphi$, $z=z$ --- обобщённые полярные, якобиан $nab\,\rho\cos^{n-1}\varphi\sin^{n-1}\varphi$. Этап 2: $\rho=r\cos^n\psi$, $z=c\,r\sin^n\psi$ --- снова обобщённые полярные, якобиан $nc\,r\cos^{n-1}\psi\sin^{n-1}\psi$. Перемножаем и подставляем $\rho=r\cos^n\psi$:
\[
\boxed{\ |J|=n^2abc\,r^2\cos^{n-1}\varphi\,\sin^{n-1}\varphi\,\cos^{2n-1}\psi\,\sin^{n-1}\psi .\ }
\]

\subsection{«Выпрямляющие» (треугольные) замены}
Если $x$ зависит только от $u$, $y$ --- от $u,v$, а $z$ --- от $u,v,w$, то матрица производных треугольная и
$J=x'_u\cdot y'_v\cdot z'_w$. Типичный приём: если $y$ меняется от $0$ до $g(x)$, кладём $y=g(x)\cdot v$, $v\in[0,1]$ --- «доля пути». Каждый такой шаг превращает криволинейную границу в плоскость $v=1$ (задачи 2 и 6).

% ===============================================================
\section{Техника однократных интегралов, которая понадобится}
% ===============================================================
\begin{enumerate}[label=\textbf{\arabic*.}]
\item \textbf{Степени косинуса по полному периоду.}
$\int_0^{2\pi}\cos t\,dt=0$;\quad
$\int_0^{2\pi}\cos^2t\,dt=\int_0^{2\pi}\frac{1+\cos2t}{2}dt=\pi$;\quad
$\int_0^{2\pi}\cos^3t\,dt=\int_0^{2\pi}(1-\sin^2t)\,d(\sin t)=\bigl[\sin t-\tfrac{\sin^3t}{3}\bigr]_0^{2\pi}=0$;\\
$\cos^4t=\Bigl(\frac{1+\cos2t}{2}\Bigr)^2=\frac38+\frac{\cos2t}{2}+\frac{\cos4t}{8}$, поэтому $\int_0^{2\pi}\cos^4t\,dt=\frac38\cdot2\pi=\frac{3\pi}{4}$.
\item \textbf{$\int\sin^m x\cos^n x\,dx$, где одна из степеней нечётная:} отщепляем одну нечётную степень и делаем замену. Например,
$\int_0^{\pi/2}\cos^3\psi\sin\psi\,d\psi=\bigl[-\tfrac{\cos^4\psi}{4}\bigr]_0^{\pi/2}=\tfrac14$;\quad
$\int_0^{\pi/2}\sin^3\varphi\cos^3\varphi\,d\varphi\overset{t=\sin\varphi}{=}\int_0^1t^3(1-t^2)\,dt=\tfrac1{12}$.
\item \textbf{Обе степени чётные:} понижаем степень:
\[\int_0^{\pi/2}\sin^2\varphi\cos^2\varphi\,d\varphi=\frac14\int_0^{\pi/2}\sin^22\varphi\,d\varphi=\frac14\int_0^{\pi/2}\frac{1-\cos4\varphi}{2}d\varphi=\frac{\pi}{16}.\]
\item \textbf{По частям:} $\int r\sin r\,dr=-r\cos r+\int\cos r\,dr=-r\cos r+\sin r+C$.
\item \textbf{Замена $t=\operatorname{tg}\varphi$:} $dt=\frac{d\varphi}{\cos^2\varphi}$, $\frac1{\cos^2\varphi}=1+t^2$ --- для интегралов вида $\int\frac{\sin^k\varphi}{\cos^{k+2m}\varphi}d\varphi$.
\end{enumerate}

% ===============================================================
\section{Алгоритм оформления: пять шагов}
% ===============================================================
Каждая задача ниже решена строго по одной схеме.

\begin{enumerate}[label=\textbf{Шаг \arabic*.},leftmargin=4em]
\item \textbf{Шапка.} Слева --- сам интеграл $\iint_{\Omega_{xy}}\dots\,dx\,dy$ или $\iiint_V\dots\,dx\,dy\,dz$. Справа --- граница ($\partial\Omega_{xy}$ или $\partial V$) с нумерацией поверхностей (1), (2), \dots\ и ограничениями знаков.
\item \textbf{Два графика рядом.} График 1 --- область (тело) в старых координатах: точки пересечения, симметрия, форма. График 2 --- та же область в новых координатах, где она «распрямилась» в прямоугольник, полосу или параллелепипед.
\item \textbf{Замена и модуль якобиана.} Система перехода в столбик, справа --- $|J|$.
\item \textbf{Пределы по второму графику.} От констант к зависимостям: сначала углы ($\varphi,\psi$), затем радиус $r$ или высота $z$.
\item \textbf{Повторный интеграл и счёт} изнутри наружу до числа.
\end{enumerate}

\begin{notebox}[Чек-лист перед ответом]
\begin{itemize}
\item Не забыт ли $|J|$? (Самая частая ошибка: $dx\,dy\ne dr\,d\varphi$, а $=r\,dr\,d\varphi$.)
\item Пределы внешнего интеграла --- числа?
\item Знак ответа разумен? (Объём --- положителен; если функция отрицательна на всей области, интеграл отрицателен.)
\item По возможности --- проверка другим способом (в декартовых координатах, по известным формулам объёма и т.\,п.). Во всех задачах ниже такая проверка сделана.
\end{itemize}
\end{notebox}
